\documentclass{article}

\usepackage{amsmath} 
\usepackage{amssymb,amsfonts,amsthm,mathtools}
\usepackage{graphicx}
\usepackage{stmaryrd}
\usepackage[inference]{semantic} %hw6

\theoremstyle{plain}
\newtheorem{theorem}{Theorem}[section]
\newtheorem{lemma}[theorem]{Lemma}
\newtheorem{corollary}[theorem]{Corollary}

\theoremstyle{definition}
\newtheorem{definition}{Definition}[section]
\newtheorem{exercise}{Exercise}[section]
\newtheorem{example}{Example}[section]

\theoremstyle{remark}
\newtheorem{remark}{Remark}[section]
\newtheorem{metatheorem}{Metatheorem}[section]
\newtheorem{axiom}{Axiom}[section]

\DeclareMathOperator{\lub}{lub}
\DeclareMathOperator{\glb}{glb}

\newcommand{\I}{\mathcal{I} } % hw6
\newcommand{\impl}{\Rightarrow } %hw6
\newcommand{\Lc}{\mathcal{L} } %hw17
\newcommand{\ato}[1]{\ensuremath{\stackrel{#1}{\to}}} %hw17
\newcommand{\atos}[1]{\ensuremath{\xrightarrow{#1}\mathrel{\vphantom{\to}^*}}} %hw17

\newcommand{\PH}{\ensuremath{\varphi}} % hw9
\newcommand{\PS}{\ensuremath{\psi}} % hw9
\newcommand{\TPR}{\ensuremath{\mathcal{T}_{PR}}} % hw9

\input{../macros}

\begin{document}
\section{Assignment 14}
For guarded BPP terms $P, P', Q$ where $P$ is a subterm of $Q$, which of the following hold?

\subsection{Language equivalence}
If $P \approx_{lin} P'$ then $Q[P] \approx_{lin} Q[P']$.
Since $P \approx_{lin} P'$ is defined as language equivalence, it must be the case that $L_T(P) = L_T(P')$.
We now consider each possible replacement for $P$ in $Q$.
\paragraph{$aP \approx_{lin} aP'$} $L_T(aP) = \{a\cdot w\mid w \in L_T(P)\} = \{a\cdot w\mid w \in L_T(P')\} = L_T(aP')$
\paragraph{$Q + P \approx_{lin} Q + P'$} 
It must be the case that $L_T(P + Q) = L_T(P) \cup L_T(Q)$ because the choice operator at the beginning of two commands forces you to choose to run one or the other; 
choose to run $P$ and you will produce any of the strings in $L_T(P)$, 
choose to run $Q$ and you will produce any of the strings in $L_T(Q)$.

Therefore $L_T(P + Q) =  L_T(P) \cup L_T(Q) = L_T(P') \cup L_T(Q) = L_T(P' + Q)$.

\paragraph{$Q \parallel P \approx_{lin} Q \parallel P'$} holds because $L_T(Q \parallel P)$ consists of all interleavings of all strings from $L_T(P)$ and $L_T(Q)$ and all possible strings where where every combination of possible $\tau$ steps is taken.
This will stay the same if $L_T(P)$ is replaced by $L_T(P')$.







\subsection{Simulation}
If $P \approx_{sim} P'$ then $Q[P] \approx_{sim} Q[P']$.

Since $P \approx_{sim} P'$ holds it must be the case that there exists some simulation relation $R_{PP'}$.
\paragraph{$aP \approx_{sim} aP'$} holds because it is possible to construct simulation relation
$$R':= \{(aP, aP')\} \cup R_{PP'}$$
$R'$ is a valid simulation relation since both $aP$ and $aP'$ can only take a single step with action $a$ and then are in a valid simulation relation.
\paragraph{$Q + P \approx_{sim} Q + P'$} holds because it is possible to construct simulation relation
$$R':= R_{QQ} \cup R_{PP'} \cup \{(Q + P, Q + P')\}$$
where $R_{QQ}$ is the simulation relation in which every state of $Q$ is paired with itself.
$R_{QQ}$ is a valid simulation relation because $\approx_{sim}$ is an equivalence relation and therefore reflexive.
$R'$ is a valid simulation relation because either the operational semantics rule applied is taking a step in $Q$ with some action, which can also be made in $Q + P'$ and then continuing in $R_{QQ}$ or taking a step in $P$, which must also be possible in $R_{PP'}$ and then continuing in that relation.

\paragraph{$Q \parallel P \approx_{sim} Q \parallel P'$} holds because it is possible to construct a simulation relation for each state $p$ in some point in process $P$ paired with every state $q$ in process $Q$, and also from every pair of states $(p,q)$ where each state is about to take synchronizing steps.

$$
R' := 
\{((p \parallel q), (p \parallel q'))\mid \forall p \in S_P, \forall (q,q') \in R_{QQ}\} \cup$$
$$\{((p \parallel q), (p' \parallel q))\mid \forall q \in S_Q, \forall (p,p') \in R_{PP'}\} \cup $$
$$\{ ((p \parallel q), (p' \parallel q')) \mid \forall a \in A, \forall p' \in S_P, p\xrightarrow{a}  p'\wedge \forall q' \in S_Q, q\xrightarrow{\bar{a}} q' \}
$$

$P'$ can take analogous steps for every $P$ so this relation will hold.









\subsection{Bisimulation}
If $P \approx_{bis} P'$ then $Q[P] \approx_{bis} Q[P']$.
Since $P \approx_{lin} P'$ is defined as language equivalence, it must be the case that $L_T(P) = L_T(P')$.
\paragraph{$aP \approx_{bis} aP'$} holds because it is possible to construct bisimulation relation
$$R':= \{(aP, aP')\} \cup R_{PP'}$$
$R'$ is a valid bisimulation relation since both $aP$ and $aP'$ can only take a single step with action $a$ and then are in a valid bisimulation relation.
\paragraph{$Q + P \approx_{bis} Q + P'$} holds because it is possible to construct simulation relation
$$R':= R_{QQ} \cup R_{PP'} \cup \{(Q + P, Q + P')\}$$
where $R_{QQ}$ is the bisimulation relation in which every state of $Q$ is paired with itself.
$R_{QQ}$ is a valid bisimulation relation because $\approx_{sim}$ is an equivalence relation and therefore reflexive.
$R'$ is a valid bisimulation relation because both (i) and (ii) hold:
\begin{itemize}
		\item[(i)] Given some $Q$ and some $P$ either the operational semantics rule applied is taking a step in $Q$ with some action, which can also be made in $Q + P'$ and then continuing in $R_{QQ}$ or taking a step in $P$ and continue in $R_{PP'}$
		\item[(ii)] Given some $Q$ and some $P"$ either the operational semantics rule applied is taking a step in $Q$ with some action, which can also be made in $Q + P$ and then continuing in $R_{QQ}$ or taking a step in $P'$ and continue in $R_{PP'}$
\end{itemize}

\paragraph{$Q \parallel P \approx_{bis} Q \parallel P'$} holds because it is possible to construct a bisimulation relation for each state $p$ in some point in process $P$ paired with every state $q$ in process $Q$, and also from every pair of states $(p,q)$ where each state is about to take synchronizing steps.

$$
R' := 
\{((p \parallel q), (p \parallel q'))\mid \forall p \in S_P, \forall (q,q') \in R_{QQ}\} \cup$$
$$\{((p \parallel q), (p' \parallel q))\mid \forall q \in S_Q, \forall (p,p') \in R_{PP'}\} \cup $$
$$\{ ((p \parallel q), (p' \parallel q')) \mid \forall a \in A, \forall p' \in S_P, p\xrightarrow{a}  p'\wedge \forall q' \in S_Q, q\xrightarrow{\bar{a}} q' \}
$$
\begin{itemize}
	\item[(i)] $P'$ can take analogous steps for every $P$ so this relation will hold.
	\item[(ii)] $P$ can take analogous steps for every $P'$ so this relation will hold.
\end{itemize}

%\begin{definition}
%\end{definition}
\end{document}
